%! Mode:: "TeX:UTF-8"
%! TEX program = xelatex
\PassOptionsToPackage{quiet}{xeCJK}
\documentclass[withoutpreface,bwprint]{cumcmthesis}
\usepackage{etoolbox}
\usepackage{ctex}
\BeforeBeginEnvironment{tabular}{\zihao{-5}}
\usepackage[numbers,sort&compress]{natbib}  % 文献管理宏包
\usepackage[framemethod=TikZ]{mdframed}  % 框架宏包
\usepackage{url}  % 网页链接宏包
\usepackage{subcaption}  % 子图宏包
\newcolumntype{C}{>{\centering\arraybackslash}X}
\newcolumntype{R}{>{\raggedleft\arraybackslash}X}
\newcolumntype{L}{>{\raggedright\arraybackslash}X}\renewcommand{\textbf}[1]{{\song\bfseries #1}}

\usepackage{tikz}
\usetikzlibrary{shapes.geometric, arrows.meta, positioning}

\tikzstyle{startstop} = [rectangle, rounded corners, minimum width=3cm, minimum height=1cm,text centered, draw=black, fill=red!20]
\tikzstyle{process} = [rectangle, minimum width=3.5cm, minimum height=1cm, text centered, draw=black, fill=blue!20]
\tikzstyle{decision} = [diamond, aspect=2, minimum width=3.5cm, minimum height=1cm, text centered, draw=black, fill=green!20]
\tikzstyle{arrow} = [thick,->,>=stealth]






















\title{微网与外部电网电力调控策略}  % 论文标题



%%%%%%%%%%%%%%%%%%%%%%%%%%%%%%%%%%%%%%%%%%%%%%%%%%%%%%%%%%%%%
%% 正文
\begin{document}
	
% 封面和目录页不显示页码
\pagenumbering{gobble}
\maketitle
\begin{abstract}


\textbf{对于问题一,}建立带周期库存约束的确定性线性规划：在母线侧能量平衡、光伏消纳与弃光、储能状态与功率边界约束下，最小化全天计划购电费用。求解得全天购电费 $35126.949$ 元、购电量 $59482.699$ kWh，较无储能基线降低 $26.8981\%$。独立逐槽复算、显式互斥 MILP 对拍与外部锚点一致，结果在冻结口径下可信。

\textbf{对于问题二,}

\textbf{对于问题三,}在问题二的因果回放框架上建立四时点确定性模型预测控制模型。附件3的整点光伏点预测经无泄漏线性展开后，在0:00、6:00、12:00和18:00依次更新；已执行时段锁定，未执行时段按相邻版本的非对称现金流重优化。334天回放总费用为15805751.77元，较问题二降低6.3912\%，紧急购电量降低32.5413\%；独立物理、账本、版本和工作簿复核19项全部通过。

\textbf{对于问题四,}

最后，



\keywords{微网；储能调度；线性规划；购电策略；光伏消纳}

\end{abstract}


%%%%%%%%%%%%%%%%%%%%%%%%%%%%%%%%%%%%%%%%%%%%%%%%%%%%%%%%%%%%% 
%\vspace{-1.4cm} 
%\tableofcontents  % 目录
%\newpage











%%%%%%%%%%%%%%%%%%%%%%%%%%%%%%%%%%%%%%%%%%%%%%%%%%%%%%%%%%%%% 
\pagenumbering{arabic}
\vspace{-1.4cm} 
\section{引言}
\subsection{问题背景}







%%%%%%%%%%%%%%%%%%%%%%%%%%%%%%%%%%%%%%%%%%%%%%%%%%%%%%%%%%%%% 
\subsection{研究意义}












%%%%%%%%%%%%%%%%%%%%%%%%%%%%%%%%%%%%%%%%%%%%%%%%%%%%%%%%%%%%% 
\subsection{问题重述}

\textbf{问题1：} 已知某天的分时电价、小区负载与光伏发电功率预测，要求微网供电不低于负载且储能设备在 0:00 与 24:00 的储电量相同，在每天 0:00 制定当天的计划购电策略，并在论文中给出指定时间段的购电量、充放电量及全天购电量与购电费。

\textbf{问题2：}  

\textbf{问题3：} 在问题二的基础上，利用每天0:00、6:00、12:00和18:00发布的未来24小时光伏预测，对尚未执行时段滚动调整购电策略；按上调1.5倍、下调0.5倍和紧急购电5倍的规则核算费用，给出全年计划购电量、调整购电量、充放电量与紧急购电量，并判断是否需要引入其他时刻的预报。

\textbf{问题4：}  

















\newpage
%%%%%%%%%%%%%%%%%%%%%%%%%%%%%%%%%%%%%%%%%%%%%%%%%%%%%%%%%%%%% 
\section{总体分析}
























\newpage
%%%%%%%%%%%%%%%%%%%%%%%%%%%%%%%%%%%%%%%%%%%%%%%%%%%%%%%%%%%%% 
\section{模型假设}

为简化问题，本文做出以下假设：

\begin{itemize}[itemindent=2em]
\item 每天 0:00 制定计划时，附件 1 中的当日电价、负荷和光伏预测均已知，且计划执行期间不再调整。
\item 按时间方案 A：附件数据、模型槽位和结果模板按位置一一对应、不平移；槽 0 对应时段 0:10-0:20，槽 143 对应次日 0:00-0:10；当天 0:00-0:10 无输入数据，不单独调度且储电量不变。
\item 微网不向外网反送电；光伏超过负荷与储能接纳能力时允许弃光。
\item 充电量和放电量均按交流母线侧计量，充、放电效率分别取 0.9。
\item 问题一采用确定性输入，不产生紧急购电、未使用合同电和调整结算费用。
\item 问题二每天 0:00 制定计划时，只能使用该时刻及之前已经观测到的历史负荷与光伏实测值；当天的真实负荷和光伏尚未发生，不得使用未来真值。
\item 问题二按计划购电量计费（含未实际使用的合同电）；供电不足部分按交易时刻电价的 5 倍紧急购电。实际执行时电池尽量按计划动作，仅在触及储电量上、下界时按需削减充电或放电量，削减量逐槽留痕。
\item 问题二跨日储电量连续传递：当日计划末槽（次日 0:00-0:10）执行后的真实储电量作为次日实际起始状态，不重置为 6000 kWh。
\item 问题二终端库存处理采用滚动 7 天视野：每天用未来 7 天的因果预测协同优化，要求视野末回到起始储电量，但只执行当天（144 槽）的计划。
\item 问题三中附件3第 $h$ 小时预报解释为发布时刻后 $h$ 小时的点预测；点间采用由一月预热期确定的分段线性展开，跨午夜仍有效的24小时预报不截断。
\item 问题三每次调整只作用于尚未执行时段；合同费用由初始计划和相邻版本调整现金流唯一确定，不再对最终计划量重复计费。
\end{itemize}

%%%%%%%%%%%%%%%%%%%%%%%%%%%%%%%%%%%%%%%%%%%%%%%%%%%%%%%%%%%%% 

\section{符号说明}









\begin{table}[H]
	\centering
	
	\begin{tabularx}{\textwidth}{CL}
		\toprule
		符号    & 说明    \\
		\midrule
		$t$ & 10 分钟时段序号，$t=0,1,\ldots,143$ \\
		$\Delta t$ & 时段长度，$\Delta t=1/6\ \mathrm{h}$ \\
		$p_t$ & 时段 $t$ 电价（元/kWh） \\
		$L_t$ & 时段 $t$ 负荷电量（kWh） \\
		$V_t$ & 时段 $t$ 可用光伏电量（kWh） \\
		$g_t$ & 时段 $t$ 外网计划购电量（kWh） \\
		$v_t$ & 时段 $t$ 实际消纳光伏电量（kWh） \\
		$q_t^{\mathrm{curt}}$ & 时段 $t$ 弃光电量（kWh） \\
		$c_t,\,d_t$ & 母线侧充电量、放电量（kWh） \\
		$E_t$ & 第 $t$ 个状态边界的储电量（kWh） \\
		$\eta_c,\,\eta_d$ & 充电效率、放电效率，均为 0.9 \\
		$g_t^{\mathrm{emg}}$ & 时段 $t$ 紧急购电量（kWh），按 $5p_t$ 计费 \\
		$g_t^{\mathrm{unused}}$ & 时段 $t$ 计划合同购电中未实际使用的部分（kWh） \\
		$\hat{L}_t,\,\hat{V}_t$ & 时段 $t$ 负荷、可用光伏的因果预测值（kWh） \\
		$\mathcal{H}_d$ & 第 $d$ 天滚动视野的预测天数（问题二取 7） \\
		$r$ & 问题三计划发布时点，$r\in\{0,6,12,18\}$ \\
		$q_t^{(r)}$ & 时段 $t$ 在发布时点 $r$ 后的合同购电计划（kWh） \\
		$a_t^{r,+},\,a_t^{r,-}$ & 相对上一计划版本的上调量、下调量（kWh） \\
$\hat V_{t\mid r}$ & 发布时点 $r$ 可获得的时段 $t$ 光伏预测（kWh） \\
$\Omega(f)$ & 正则化项 \\

		\bottomrule
	\end{tabularx}
	
	\label{tab:符号说明}
\end{table}











\newpage
%%%%%%%%%%%%%%%%%%%%%%%%%%%%%%%%%%%%%%%%%%%%%%%%%%%%%%%%%%%%% 
\section{问题一的模型的建立和求解}
\subsection{具体分析}

问题一给出了一个典型日内完整的分时电价、小区负荷和光伏发电预测，因此不需要再建立预测模型，而是要在满足负荷与储能约束的前提下，确定 144 个 10 分钟时段的计划购电量及储能充放电量。该问题的核心是利用低价电和光伏富余给电池充电，在高价时段放电，同时避免通过期末耗空电池获得虚假的低成本。由此可将问题建立为带周期库存约束的确定性线性规划。

\subsection{模型准备}

每天 0:00 制定计划时，附件 1 中的当日电价、负荷和光伏预测均已知，且计划执行期间不再调整。按冻结的时间方案 A：附件数据、模型槽位和结果模板按位置一一对应、不平移，槽 0 对应时段 \texttt{0:10-0:20}，槽 143 对应次日 \texttt{0:00-0:10}；附件缺少当天 \texttt{0:00-0:10} 的输入，该段不单独调度且储电量不变。微网不向外网反送电；光伏超过负荷与储能接纳能力时允许弃光。充电量和放电量均按交流母线侧计量，充、放电效率分别取 0.9。问题一采用确定性输入，不产生紧急购电、未使用合同电和调整结算费用。

令 $t=0,1,\ldots,143$ 表示 10 分钟时段，$\Delta t=1/6\ \mathrm{h}$。附件 1 中负荷和光伏为功率量，对每个时段分别计算电量
\begin{equation}
L_t=P_t^{\mathrm{load}}\Delta t,\qquad
V_t=P_t^{\mathrm{pv}}\Delta t. \label{eq:q1_energy}
\end{equation}
读取时要求价格、负荷和光伏恰有 144 个数值，无缺失和重复；时间映射统一由槽位编号完成，不使用模糊的时间字符串连接。

\subsection{模型建立}

每个时段的母线侧能量平衡为
\begin{equation}
g_t+v_t+d_t=L_t+c_t, \label{eq:q1_balance}
\end{equation}
其中 $g_t$ 为外网计划购电量，$v_t$ 为实际消纳光伏电量，$c_t,d_t$ 为母线侧充电量、放电量。光伏消纳及弃光满足
\begin{equation}
0\le v_t\le V_t,\qquad q_t^{\mathrm{curt}}=V_t-v_t\ge0. \label{eq:q1_pv}
\end{equation}
储能状态按充放电效率递推：
\begin{equation}
E_{t+1}=E_t+\eta_c c_t-\frac{d_t}{\eta_d},\qquad
\eta_c=\eta_d=0.9. \label{eq:q1_soc}
\end{equation}
容量、功率和购电约束为
\begin{equation}
1200\le E_t\le10800,\qquad
0\le c_t,d_t\le5000\Delta t=833.333,\qquad
g_t\ge0. \label{eq:q1_bounds}
\end{equation}
按照冻结的方案 A，完整执行 144 个模板槽后回到初始库存：
\begin{equation}
E_0=E_{144}=6000. \label{eq:q1_terminal}
\end{equation}
主目标为最小化全天正常购电费用：
\begin{equation}
\min C=\sum_{t=0}^{143}p_tg_t. \label{eq:q1_obj}
\end{equation}
由于线性规划可能存在多个同费用解，先求得最优费用 $C^*$，再在 $C\le C^*+10^{-7}$ 的条件下最小化 $\sum_t(c_t+d_t)$。该字典序目标只消除无意义的充放电循环，不改变经济目标。

\subsection{模型求解}

首先用 HiGHS dual simplex 求解主模型，并用 HiGHS IPM 对拍目标值；随后将购电、光伏消纳、充放电和储电量送入独立执行器。核算器不调用优化器目标函数，而是重新检查式 \eqref{eq:q1_energy} 至式 \eqref{eq:q1_terminal}，逐槽计算 $p_tg_t$，最后把审计通过的结果写入官方模板并重新打开回读。论文中的费用均取核算器结果，而不是直接复制求解器日志。

表 \ref{tab:q1_grid} 给出题面指定时段的购电量。完整 144 槽结果保存在正式结果文件 \texttt{result1.xlsx} 中。

\begin{table}[H]
	\centering
	\caption{问题一指定时间段购电量及全天购电量、购电费}
	\begin{tabularx}{\textwidth}{lrlrlr}
		\toprule
		时间段 & 购电量（kWh） & 时间段 & 购电量（kWh） & 时间段 & 购电量（kWh） \\
		\midrule
		10:00-10:10 & 0.000 & 12:00-12:10 & 486.403 & 14:00-14:10 & 0.000 \\
		16:00-16:10 & 394.931 & 18:00-18:10 & 636.983 & 20:00-20:10 & 0.000 \\
		\bottomrule
	\end{tabularx}
	\label{tab:q1_grid}
\end{table}

全天计划购电量为 $59482.699$ kWh，购电费为 $35126.949$ 元。光伏消纳量为 $55482.836$ kWh，弃光量为 0；累计充电量和放电量分别为 $20740.666$ kWh 与 $16799.940$ kWh。

表 \ref{tab:q1_cd} 按题面六个 4 小时块汇总充放电量。为与时间方案和结果模板一致，第 $i$ 个 24 槽按位置归入对应块，不另做十分钟平移。

\begin{table}[H]
	\centering
	\caption{问题一储能设备分块充放电量与 0:00、24:00 储电量}
	\begin{tabularx}{\textwidth}{lrlrlr}
		\toprule
		时间段 & 充电量（kWh） & 放电量（kWh） & 时间段 & 充电量（kWh） & 放电量（kWh） \\
		\midrule
		0:00-4:00 & 4500.000 & 0.000 & 4:00-8:00 & 833.333 & 6365.841 \\
		8:00-12:00 & 4787.964 & 1702.997 & 12:00-16:00 & 5286.035 & 91.101 \\
		16:00-20:00 & 0.000 & 5780.132 & 20:00-24:00 & 5333.333 & 2859.868 \\
		\bottomrule
	\end{tabularx}
	\label{tab:q1_cd}
\end{table}

方案中储电量最小值和最大值分别为 $1200$ kWh 和 $10800$ kWh，均位于允许范围内；按题面表 2 口径，0:00 与 24:00 储电量均为 $6000$ kWh。与不使用储能的可行基线相比，基线费用为 $48052.047$ 元，本方案减少 $12925.098$ 元，降幅为 $26.8981\%$。这表明储能对分时电价和光伏时移具有明显经济价值。

问题一的模型检验内容见“模型的分析与检验”一节：独立逐槽复算最大残差不超过 $10^{-12}$ 量级，显式互斥 MILP 对拍费用差 $7.28\times10^{-12}$ 元，效率口径敏感性约 $\pm3.8\%$。敏感性细节同样在该节给出，此处不再重复。

综上，问题一建立了四问共用的物理和核算内核。确定性 LP 在满足供需平衡、储电量、功率和周期库存约束的情况下，将购电费用降至 $35126.949$ 元，较无储能方案降低 $26.8981\%$。当前主要局限来自题面时间标签和效率语义，而非求解器或结果文件实现。

\newpage
%%%%%%%%%%%%%%%%%%%%%%%%%%%%%%%%%%%%%%%%%%%%%%%%%%%%%%%%%%%%% 
\section{问题二的模型的建立和求解}
\subsection{具体分析}

问题二要求：在每天电价相同、负荷与光伏随时间变化、供电不可低于负荷（不足则按交易时刻电价的 5 倍紧急购电）的条件下，每天 0:00 制定当天的 144 段购电与充放电计划，并给出 2025 年 2 月 1 日至 12 月 31 日的完整策略。与问题一最大的不同是：制定计划时当天的真实负荷与光伏尚未发生，计划只能依据截至 0:00 已经观测到的历史数据，因此预测误差会直接触发 5 倍惩罚。核心难点有三：一是严格因果的负荷与光伏预测；二是跨日储电量连续传递（不能每天重置为 6000 kWh）；三是终端库存口径——每天只优化当天会把期末库存耗空换取虚假低价，需要合适的跨日前瞻机制。

\subsection{模型准备}

\textbf{信息集。} 每天 0:00 的决策只能使用该时刻及之前已经完整观测到的实际负荷与光伏。特别地，按时间方案 A，前一天计划的最后一槽位于次日 0:00 至 0:10，要等到 0:10 才观测完整，因此 0:00 制定计划时不能读取它。所有历史输入满足“观测时刻 $\le$ 决策时刻”。

\textbf{跨日零点衔接。} 0:00 至 0:10 继续执行前一天已经制定的最后一槽计划；新一天的计划在 0:00 制定、从 0:10 开始执行。因此 0:00 制定新计划时，只用当时已知的实际储电量与旧计划已安排的充放电动作，估计 0:10 的储电量作为新计划的起点；0:10 得到真实储电量后，估计偏差由储能边界的按需削减吸收，次日继续传递真实储电量。

\textbf{数据与冷启动。} 负荷取附件 2 的“小区负载”实测，光伏取附件 2 的“光伏发电实际功率”实测（本机官方包该表缺失，按既有裁决使用与官方一致口径的替补光伏表）。由于 2 月 1 日需要由一月的实际执行结果推出起始储电量，1 月按同样的执行器顺序回放：1 月 1 日只有附件 1 的典型日先验，用其作为负荷与光伏预测；1 月 2 日至 7 日优先使用决策时已经完整观测到的同槽历史均值，尚无观测的槽继续回退典型日先验；1 月 8 日起启用历史预测基线。2 月 1 日不重置，直接继承 1 月 31 日的实际末储电量。

\textbf{执行规则。} 电池尽量按 0:00 的计划动作，不因实际值偏离而实时重调；供电缺口全部走紧急购电。当实际负荷低于预测导致计划放电被削减、进而使后续计划充电会超过 10800 kWh 上限时，执行器仅为保持物理边界而削减必要的充电量。所有削减逐槽记为计划值、实际值与缺口量，供独立核算器复查。实际执行遵循“先弃光、再形成未使用合同电、最后削减计划放电”的优先顺序。

\textbf{费用账本。} 设时段 $t$ 的计划合同购电量为 $g_t$、实际进入母线的合同电为 $g_t^{\mathrm{deliv}}$、紧急购电量为 $g_t^{\mathrm{emg}}$，则费用为
\begin{equation}
C=\sum_{t} p_t\,g_t \;+\; \sum_{t} 5\,p_t\,g_t^{\mathrm{emg}},
\end{equation}
其中未实际使用的合同电 $g_t^{\mathrm{unused}}=g_t-g_t^{\mathrm{deliv}}$ 不进入母线能量平衡，也不额外收费（已含在正常合同电费中）。母线能量平衡为
\begin{equation}
g_t^{\mathrm{deliv}}+g_t^{\mathrm{emg}}+v_t+d_t = L_t + c_t .
\end{equation}

\subsection{模型建立}

\textbf{预测基线。} 采用七天前同一墙钟时刻的实际值作为当天负荷与光伏的点预测（记 D-7）：
\begin{equation}
\hat{L}_t^{(d)}=L_t^{(d-7)},\qquad \hat{V}_t^{(d)}=V_t^{(d-7)},
\end{equation}
其优点是完全因果、无需拟合、能保留完整的周内规律。

\textbf{滚动 7 天视野的日计划。} 每天 0:00 以未来 7 天的因果预测协同优化，要求视野末储电量回到起始储电量，只执行当天 144 槽的计划，次日滚动重算。目标为
\begin{equation}
\min \sum_{k=0}^{7\times144-1} p_{k}\,\hat{g}_{k},
\end{equation}
约束为母线能量平衡、光伏消纳与弃光、充放电功率上限 $\bar c=\bar d=5000\Delta t=\tfrac{5000}{6}$ kWh、储电量边界 $1200\le E\le 10800$ 及视野周期末端 $E_{\mathrm{end}}=E_0$。该口径用“未来几天要回来”替代对期末库存的人为定价，既避免每日耗空电池，也避免把库存长期钉在上限而丧失调节能力。

\subsection{模型求解}

以 D-7 点预测与滚动 7 天视野对 2025 年 1 月 1 日至 12 月 31 日做顺序回放（1 月为预热），2 月 1 日至 12 月 31 日共 334 天的结果由独立核算器复算：全天总购电费 \textbf{16884897.80 元}，其中正常合同购电费 12190821.03 元、紧急购电费 4694076.77 元，紧急购电量 1199589.73 kWh，未使用合同电 315194.51 kWh，不可行天数为 0，全部逐槽能量平衡与储电量连续检查通过。年末实际储电量 8824.35 kWh，全年日末储电量中位数约 8663 kWh，既未长期贴上限也未贴下限。

\textbf{预测方法的费用比较。} 在同一滚动 7 天视野下，七天前同刻全年总费用 16884897.80 元，同星期日四周衰减平均为 17296322.75 元，前者低约 411425 元。虽然同星期日衰减平均在 2 至 4 月的短窗口内误差更小，但长程上 D-7 保留周内规律、误差不随季节漂移累积，最终费用更低——这说明预测误差大小不能单独决定策略优劣。

\textbf{保守程度的经济性。} 我们另备了将计划量按历史预测残差的分位数（0.7/0.8/0.9）上调的保守候选。直觉上 5 倍紧急电价应使“报童临界分位”落在 0.8 附近，但费用比较显示点预测反而最省：在 2 至 4 月窗口内，点预测总费用约 357 万元，0.7、0.8、0.9 分位分别约 390、400、415 万元，保守程度越高越贵。原因是储能本身已经跨时段对冲了预测误差，在购电侧再按分位加保守裕度属于重复对冲，只增加正常购电费。由此确认主口径为点预测。

最终第二问的完整购电策略写入按官方模板补全的 \texttt{result2.xlsm}：计划购电量表 334 天 $\times$ 144 槽全部填写，另按题面补齐“充放电量”（分块充放电量与 0:00、24:00 储电量）和“紧急购电量”（合并后的连续时段与电量）工作表；官方原有工作表、顺序、隐藏状态与 VBA 容器保持不变，保存后重新打开回读核验通过。




























\newpage
%%%%%%%%%%%%%%%%%%%%%%%%%%%%%%%%%%%%%%%%%%%%%%%%%%%%%%%%%%%%% 
\section{问题三的模型的建立和求解}
\subsection{具体分析}

问题三的新增信息不是当天光伏真值，而是日内逐次发布的未来24小时光伏预报。因此，本问的核心不是重新建立一套储能模型，而是在问题二的物理执行器上加入“预测发布--计划改版--按版结算”的闭环。预报更新可以减少预测偏差引起的紧急购电，但频繁改版又会产生上调溢价和下调冲回，故是否需要6:00、12:00、18:00预报必须由实际总费用、紧急购电量和调整规模共同判断，不能只比较预测误差。

\subsection{模型准备}

\textbf{信息集与时间映射。} 0:00先形成初始计划，6:00、12:00和18:00分别加入当时刚发布的附件3光伏预报；任一时刻只使用发布时间不晚于该时刻的数据，已经执行的槽位永久锁定。按照时间方案A，三个调整边界分别对应槽位35、71和107，故新版本从 $[6{:}00,6{:}10)$、$[12{:}00,12{:}10)$ 和 $[18{:}00,18{:}10)$ 起生效。附件3第 $h$ 列保留为 $\mathrm{valid\_time}=\mathrm{issue\_time}+h$ 的点预测，不把首小时左移。

\textbf{点预测展开。} 对相邻整点预报作分段线性插值；发布时刻的左端锚点只取此前已经发布的批次，2025年1月1日0:00无历史锚点的首小时回退附件1先验。用不进入正式评价期的一月选择展开方法：4429个可比较槽上，线性展开的MAE、RMSE分别为61.93 kW、133.09 kW，前值保持分别为242.04 kW、463.22 kW，故采用线性展开。当前批次跨午夜仍有效的24小时预测完整保留，超过覆盖范围后才回退严格因果的D-7光伏基线；负荷预测仍沿用问题二的D-7方法。

\textbf{滚动视野。} 每次求解均保留问题二的7天视野和视野末库存回归约束，只执行到下一发布时点。当前批次对未来24小时提供光伏信息，更远时段用D-7补齐；这种处理既利用新信息，又避免短视地在更新边界耗空电池。0:00跨日衔接仍按D009执行：先依据旧计划末槽动作估计新计划起始库存，再执行旧日最后一槽，实际库存连续传给下一日。

\subsection{模型建立}

设 $q_t^{(r^-)}$ 为调整前最近一版合同量，$q_t^{(r)}$ 为发布时点 $r$ 的新版合同量。只对 $t\ge s_r$ 的未执行时段定义非负调整量：
\begin{equation}
q_t^{(r)}=q_t^{(r^-)}+a_t^{r,+}-a_t^{r,-},\qquad a_t^{r,+},a_t^{r,-}\ge0.
\end{equation}
0:00初始计划沿用问题一的供需平衡、光伏、功率和SOC约束，在7天预测视野内最小化普通购电费。日内更新时，对当日已经签约但尚未执行的部分，目标函数改为相邻版本的调整现金流；对视野中尚未签约的后续日期仍按普通购电费计价：
\begin{equation}
\min\left\{
\sum_{t=s_r}^{143}\left(1.5p_ta_t^{r,+}-0.5p_ta_t^{r,-}\right)
+\sum_{t\in\mathcal H_r^{\mathrm{future}}}p_tg_t
\right\}.
\end{equation}
两部分共享同一SOC状态链和视野末约束。求得主目标最优值后，再以充放电与调整总量最小作词典序第二目标，消除同价时段无意义的往返动作。

实际执行时，合同电、电池和光伏仍遵循问题二的统一执行规则；预测不足形成的缺口由5倍电价紧急购电补齐。全年总费用由版本账本唯一核算：
\begin{equation}
C_3=\sum_t p_tq_t^{(0)}
+\sum_{r\in\{6,12,18\}}\sum_{t\ge s_r}
\left(1.5p_ta_t^{r,+}-0.5p_ta_t^{r,-}\right)
+\sum_t5p_tg_t^{\mathrm{emg}} .
\end{equation}
其中下调项是按冻结口径对被替换合同量的0.5倍冲回，不是额外罚款；最终合同量 $q_t^{\mathrm{final}}$ 不再重复收取普通购电费。

\subsection{模型求解}

采用HiGHS对每个发布时点的稀疏线性规划顺序求解。2025年1月用于预测与库存预热，正式评价期为2月1日至12月31日共334天。四时点主方案的费用与电量分解见表~\ref{tab:q3_main_result}。

\begin{table}[H]
\centering
\caption{问题三四时点滚动策略的全年结果}
\begin{tabular}{lrr}
\toprule
类别 & 数值 & 单位 \\
\midrule
初始合同购电费 & 12317106.21 & 元 \\
上调购电费 & 600179.55 & 元 \\
下调冲回 & 244787.55 & 元 \\
净调整现金流 & 355392.01 & 元 \\
紧急购电费 & 3133253.56 & 元 \\
总购电费 & \textbf{15805751.77} & 元 \\
上调量/下调量 & 618742.57/754150.54 & kWh \\
紧急购电量/未用合同电量 & 809227.14/304721.36 & kWh \\
\bottomrule
\end{tabular}
\label{tab:q3_main_result}
\end{table}

相较问题二仅在0:00制定计划的正式策略，问题三总费用减少1079146.02元，降幅为6.3912\%；紧急购电量减少390362.60 kWh，降幅为32.5413\%。这表明在冻结的非对称结算规则下，日内更新带来的缺口规避收益超过净调整成本。334天均可行，SOC始终位于1200--10800 kWh，年末SOC为8674.35 kWh。

%% Q3_UPDATE_ABLATION_BEGIN
\textbf{更新时点消融。} 保持数据、执行器、结算器和预测展开完全一致，依次加入6:00、12:00和18:00批次，结果见表~\ref{tab:q3_ablation}。6:00使总费用下降341695.91元、紧急购电量下降159601.56 kWh；12:00在此基础上再降低203250.02元和70515.50 kWh。18:00相较截至12:00的方案反而轻微增加2.90元和0.27 kWh，变化均低于0.001\%。因此，6:00与12:00预报具有明确经济价值，18:00在当前数据上不具实质边际收益；没有证据支持再引入更密集的其他发布时间。由于附件未提供其他发布时间的预报，本文不对其作无法验证的收益推断。

\begin{table}[H]
\centering
\caption{日内光伏预报更新时点的逐级消融}
\begin{tabular}{lrrrr}
\toprule
可用发布时点 & 总费用/元 & 净调整费/元 & 紧急购电量/kWh & 不可行天数 \\
\midrule
00 & 16350694.80 & 0.00 & 1039343.92 & 0 \\
00+06 & 16008998.90 & 271160.45 & 879742.36 & 0 \\
00+06+12 & \textbf{15805748.88} & 355391.71 & \textbf{809226.86} & 0 \\
00+06+12+18 & 15805751.77 & 355392.01 & 809227.14 & 0 \\
\bottomrule
\end{tabular}
\label{tab:q3_ablation}
\end{table}

\begin{figure}[H]
\centering
\includegraphics[width=0.95\textwidth]{figures/q3_cost_and_updates.png}
\caption{问题二、问题三月度费用及更新时点消融}
\label{fig:q3_cost_updates}
\end{figure}
%% Q3_UPDATE_ABLATION_END

\textbf{独立验证与结果文件。} 独立脚本不调用优化器，直接由48096条逐槽记录和121242条计划版本复算。19项检查全部通过：母线平衡和SOC递推最大残差分别为 $7.96\times10^{-13}$ kWh 和 $5.46\times10^{-12}$ kWh，跨日SOC断裂为0，时间轴最大断裂为0秒；版本顺序、目标槽定价、逐日费用与最终版本映射均一致。正式工作簿逐槽购电量最大回读误差为 $4.55\times10^{-13}$ kWh，6407个合并紧急购电区间的总量最大误差为 $1.82\times10^{-12}$ kWh。最终策略已写入 \texttt{result3.xlsm}，包括334天初始购电量、调整购电量、充放电量及紧急购电量。

上述结论是在固定日内电价、D-7负荷预测、线性光伏展开和D007版本结算口径下得到的样本外回放结果。每次LP在给定信息集下最优，但不把它表述为不确定环境下的全年全局最优；更新时点的必要性由下述逐级消融进一步判断。
















\newpage
%%%%%%%%%%%%%%%%%%%%%%%%%%%%%%%%%%%%%%%%%%%%%%%%%%%%%%%%%%%%% 
\section{问题四的模型的建立和求解}
\subsection{具体分析}

\subsection{模型准备}

\subsection{模型建立}

\subsection{模型求解}











\newpage
%%%%%%%%%%%%%%%%%%%%%%%%%%%%%%%%%%%%%%%%%%%%%%%%%%%%%%%%%%%%%
\section{模型的分析与检验}

\subsection{灵敏度分析}

初始储电量在 1200/6000/10800 kWh 间变化时，费用最大增加 $0.2143\%$；最大充放电功率在 4000/5000/6000 kW 间变化时，费用变化范围为 $-0.1243\%$ 至 $+0.3442\%$。这两类参数不会改变储能显著降低费用的结论。

效率是最敏感的参数：把充、放电效率由 0.9 相对降低 $5\%$ 至 0.855，费用增加 $3.7606\%$；提高至 0.945，费用降低 $3.4978\%$。若把“90\%”解释为往返效率并令 $\eta_c=\eta_d=\sqrt{0.9}$，费用为 $33801.496$ 元，比主口径低 $3.7733\%$。因此效率口径必须在全文保持一致，并把另一解释保留为敏感性结果。

对负荷和光伏施加 $1\%$、$3\%$ 和 $5\%$ 的有界 AR(1) 时间相关扰动，每级运行 20 次。60 次模型均可行且无同时充放电；相对无储能的节省率始终位于 $23.97\%$–$29.80\%$。这说明小幅输入变化会改变精确费用，但不会推翻储能调度具有经济优势的核心结论。

\begin{figure}[H]
	\centering
	\includegraphics[width=0.85\textwidth]{figures/q1_sensitivity_robustness.png}
	\caption{参数敏感性与输入稳健性}
	\label{fig:q1_sens}
\end{figure}

\subsection{误差分析}

独立验证代码不调用生产 LP、执行器和核算器，而是直接从附件 1 与最终 CSV 逐槽复算。144 槽无缺失、重复或非有限值；负荷和光伏的单位换算最大误差分别为 $1.14\times10^{-13}$ kWh 和 $2.27\times10^{-13}$ kWh；能量平衡最大残差为 $2.27\times10^{-13}$ kWh，储电量转移最大误差为 $1.82\times10^{-12}$ kWh。储电量、功率和购电语义违反次数均为 0，CSV、论文汇总表和正式工作簿的购电量最大差为 $6.82\times10^{-13}$ kWh。

进一步建立显式保留 145 个储电量状态、用二进制变量强制充放电互斥的 MILP。该模型费用为 $35126.948589389634$ 元，与主 LP 的绝对差仅 $7.28\times10^{-12}$ 元；逐槽购电 Pearson 相关系数为 $0.995576$，且两个模型均无同时充放电。最大逐槽差 $333.333$ kWh 出现在价格相同的两个时段，属于同费用多重最优解，不影响总费用和最终储电量。因此，主结果同时得到逐槽物理账本和替代优化模型支持。

\begin{figure}[H]
	\centering
	\includegraphics[width=0.85\textwidth]{figures/q1_alternative_validation.png}
	\caption{主模型与替代模型验证}
	\label{fig:q1_alt}
\end{figure}

问题一是确定性优化，不包含拟合或预测参数，因此传统随机 K 折验证并不适用。输入扰动仅作为局部压力测试，不能替代问题二至问题四的全年顺序回测。

\newpage
%%%%%%%%%%%%%%%%%%%%%%%%%%%%%%%%%%%%%%%%%%%%%%%%%%%%%%%%%%%%%
\section{模型的评价、改进与推广}

\subsection{模型的优点}
\begin{itemize}[itemindent=2em]
\item 优点1
\item 优点2
\item 优点3
\end{itemize}

\subsection{模型的缺点}
\begin{itemize}[itemindent=2em]
\item 缺点1
\item 缺点2
\end{itemize}




\subsection{模型的改进}






\subsection{模型的推广}






































\newpage
%%%%%%%%%%%%%%%%%%%%%%%%%%%%%%%%%%%%%%%%%%%%%%%%%%%%%%%%%%%%%
%% 参考文献
\nocite{*}
\bibliographystyle{gbt7714-numerical}  % 引用格式
\bibliography{ref.bib}  % bib源














\newpage
%%%%%%%%%%%%%%%%%%%%%%%%%%%%%%%%%%%%%%%%%%%%%%%%%%%%%%%%%%%%%
%% 附录
\begin{appendices}
	
	
	
	
\section{文件列表}
\begin{table}[H]
\centering
\begin{tabularx}{\textwidth}{LL}
\toprule
文件名   & 功能描述 \\
\midrule
slot\_adapter.py & 144 槽位适配器 \\
lp\_kernel.py & 确定性 LP 内核与双引擎求解 \\
executor.py & 执行器与计划版本账本 \\
accountant.py & 独立核算器 \\
template\_io.py & 官方模板原位填写与回读 \\
data\_io.py & 附件读取与防泄漏接口 \\
generate\_q1\_result.py & Q1 正式结果生成与审计 \\
verify\_q1\_result.py & Q1 结果独立交叉验证 \\
q3\_forecast.py & 附件3点预测的因果10分钟展开 \\
q3\_optimizer.py & Q3初始计划与日内调整的稀疏LP \\
q3\_replay.py & 四时点顺序回放与合同版本记录 \\
generate\_q3\_result.py & Q3全年正式结果生成 \\
validate\_q3\_result.py & Q3物理、账本与工作簿独立验证 \\
\bottomrule
\end{tabularx}
\label{tab:文件列表}
\end{table}









%\section{代码}
%\noindent 
%%%%%%%%%%%%%%%%%%%%%%%%%%%%%%%%%%%%%%%%%%%%
%数据预处理.py
%\lstinputlisting[language=python]{code/1.数据预处理.py}
%%%%%%%%%%%%%%%%%%%%%%%%%%%%%%%%%%%%%%%%%%%%%
%
%%%%%%%%%%%%%%%%%%%%%%%%%%%%%%%%%%%%%%%%%%%%
%q2.py
%\lstinputlisting[language=python]{code/q2.py}
%%%%%%%%%%%%%%%%%%%%%%%%%%%%%%%%%%%%%%%%%%%%
%
%
%q3.c
%\lstinputlisting[language=c]{code/q3.c}
%
%
%q4.cpp
%\lstinputlisting[language=c++]{code/q4.cpp}

\end{appendices}
\end{document}
